% ============================================================
% Chapter 6: Experimental Design
%   Section: Selection and Justification of the Evaluation Dataset
%     Subsection: X-Fact as an Open-Retrieval Benchmark
% Included from chapters/06-experimental-design/02-evaluation-dataset/evaluation-dataset.tex
% ============================================================

\subsection{X-Fact as an Open-Retrieval Benchmark}
\label{subsec:xfact}

As an alternative, \textbf{X-Fact} \cite{gupta2021xfact} was
selected: the largest publicly available multilingual dataset for
the verification of naturally occurring, real-world claims. X-Fact
contains 31,189 claims across 25 languages --- including Spanish ---
drawn from professional fact-checking organisations and labelled by
expert fact-checkers using a fine-grained veracity scale (not limited
to FEVER's supported/refuted/not-enough-info triad).

The decisive structural difference with respect to FEVER is that
X-Fact's claims are \textbf{real-world claims} made in public
discourse (statements by public figures, news headlines, viral
posts), not synthetic sentences derived from Wikipedia, and their
verification --- both by the original human fact-checkers and by the
baseline models proposed by the authors --- relies on evidence
retrieved via search engines over the open web, not against a
predefined closed corpus. The authors themselves report that their
best baseline model, which combines the claim text with metadata and
evidence retrieved through a search engine, achieves an $F_1$ of
approximately 40\%, positioning X-Fact as a challenging benchmark
even for systems that already operate under the same open-retrieval
paradigm as the method proposed in this work.

This alignment of paradigms --- open retrieval in the evaluation
dataset and open retrieval in the evaluated system --- is what makes
X-Fact a methodologically coherent choice: it allows the full
pipeline (Selection, Retrieval, Ranking, and LLM check jointly) to be
measured, not only its final reasoning stage, and situates the
results obtained within a framework comparable to the state of the
art reported by Gupta and Srikumar \cite{gupta2021xfact} for the
Spanish-language subset.
